%% CV tiếng Việt — Võ Ngọc Khả Ái. Bản tối giản (ATS-style). Biên dịch: ./build.sh cv-classic.tex
%% Nguồn: Báo cáo thực tập Sơn Hải Vân + danh sách đề tài KLTN IUH.
%% Chỗ [trong ngoặc vuông] là thông tin còn thiếu — bạn tự điền.
\documentclass[vietnamese]{cvclassic}

\cvdensity{0.95}

\cvname{Võ Ngọc}{Khả Ái}
\cvtitle{Kỹ sư Công nghệ Kỹ thuật Hóa học --- QA/QC \& R\&D}

\cvcontact{}{khaai.iuh@gmail.com}{mailto:khaai.iuh@gmail.com}
\cvcontact{}{0822 909 029}{tel:+84822909029}
\cvcontact{}{TP. Hồ Chí Minh}{}

\cvfooter{Võ Ngọc Khả Ái · CV · 2026}

\begin{document}
\makecvheader

\section{Giới thiệu}
\cvsummary{%
Sinh viên năm cuối ngành \textbf{Công nghệ Kỹ thuật Hóa học}, Trường Đại học
Công nghiệp TP.HCM (IUH), định hướng chuyên sâu \textbf{hóa hữu cơ}. Đã thực tập
tại \textbf{Công ty TNHH Sơn Hải Vân}, tiếp cận trọn chuỗi sản xuất sơn từ nguyên
liệu (nhựa alkyd, acrylic, epoxy, polyester; dung môi; bột màu --- bột độn) qua đánh
paste, nghiền mịn, pha loãng đến kiểm soát chất lượng đầu ra; trực tiếp vận hành và
đọc kết quả trên các thiết bị đo cơ lý tính màng sơn (độ nhớt, độ phủ, độ mịn, tỷ
trọng, độ bóng, độ bám dính) theo quy trình và biểu mẫu chuẩn hóa. Song song,
thực hiện khóa luận về \textbf{vật liệu composite giả da từ lục bình và phụ phẩm
nông nghiệp}. Sử dụng được các thiết bị đo lường kỹ thuật và phân tích cơ bản tại
phòng thí nghiệm nhà trường. Làm việc cẩn thận, có trách nhiệm với số liệu kiểm nghiệm và quen nhịp
trực ca của nhà máy. Mong muốn làm \textbf{nhân viên QA/QC hoặc R\&D}
tại doanh nghiệp sơn, chất phủ, vật liệu hoặc hóa chất công nghiệp.}

\section{Kinh nghiệm thực tế}
\cventry{Thực tập sinh --- Bộ phận Kỹ thuật (QA/QC)}{Công ty TNHH Sơn Hải Vân --- CBHD: Lê Quốc Thời, Trưởng phòng Kỹ thuật}{20/04/2026 -- 20/06/2026}{Tây Ninh}
\begin{cvitems}
  \item \textbf{Kiểm soát chất lượng theo tiêu chuẩn}: thực hiện quy trình lấy mẫu và kiểm tra ba chặng \textbf{IQC --- IPQC --- FQC} trên dây chuyền mẻ công nghiệp \textbf{400--500 kg}; ghi kết quả vào biểu mẫu chuẩn hóa của hệ thống \textbf{ISO 9001:2015}.
  \item \textbf{Trực tiếp vận hành thiết bị đo cơ lý tính màng sơn}: độ nhớt Stormer (KU), độ phủ Pfund Cryptometer, độ mịn Hegman, tỷ trọng, độ bóng (đo ở 20°, 60°, 85°) và độ bám dính --- kiểm ở \textbf{25\,±\,0,5\,°C} do độ nhớt KU rất nhạy nhiệt, đối chiếu dung sai kỹ thuật (ví dụ tỷ trọng \textbf{1,25\,±\,0,02\,g/cm³}, sơn lót alkyd \textbf{\textasciitilde70 KU}) để kết luận đạt/không đạt.
  \item \textbf{Nắm quy trình thử độ bền dài hạn} tại phòng thử nghiệm của nhà máy: lão hóa nhanh \textbf{QUV}, thử uốn trục hình trụ (TCVN 2099/ISO 1519), thử va đập (TCVN 2100-2/ISO 6272-2) và \textbf{phun sương muối NaCl 5\%, 35\,°C, 1000 giờ} (ISO 7253) --- gồm cách chuẩn bị mẫu và đánh giá độ gỉ, phồng rộp (ISO 4628-2), độ lan rỉ từ vết rạch chữ X (ASTM D1654).
  \item \textbf{So màu \& ngoại quan}: dùng tủ soi màu và máy đo màu đối chiếu mẫu với chuẩn khách hàng --- \textbf{phân biệt màu tốt, không mù màu}.
  \item \textbf{Điều chỉnh mẻ theo chỉ tiêu}: nắm nguyên tắc bổ sung dung môi từng bước để kéo độ nhớt về ngưỡng kỹ thuật, tuyệt đối không hạ quá mức vì mẻ khó cứu ngược.
  \item \textbf{Nắm dây chuyền sản xuất}: máy đánh paste (khuấy phân tán tốc độ cao), máy nghiền rổ, máy pha loãng thủy lực --- hiểu ảnh hưởng của nhiệt độ, thời gian nghiền, tốc độ bơm nhựa tới chất lượng mẻ.
  \item \textbf{Tìm hiểu cách xử lý sự cố mẻ}: theo sát cách kỹ sư nhà máy truy nguyên nhân gốc rễ khi mẻ sơn sai chỉ tiêu (sốc nhựa do bơm quá nhanh, keo tụ bột màu, lệch tỷ trọng/độ nhớt/độ mịn) và điều chỉnh để thu hồi mẻ thay vì loại bỏ.
  \item \textbf{An toàn --- môi trường (HSE)}: tuân thủ PCCC và PPE, tra cứu MSDS, phân loại và lưu chứa dung môi thải, chất thải rắn nguy hại theo quy trình nhà máy.
\end{cvitems}

\section{Học vấn}
\cventry{Kỹ sư Công nghệ Kỹ thuật Hóa học}{Trường Đại học Công nghiệp TP.HCM (IUH) --- Khoa Công nghệ Hóa học}{2022 -- 2026}{TP. Hồ Chí Minh}
\begin{cvitems}
  \item GPA \textbf{2,8}/4.0.
  \item \textbf{Khóa luận tốt nghiệp}: \textit{Nghiên cứu các yếu tố ảnh hưởng đến quy trình sản xuất vật liệu composite, vật liệu giả da từ lục bình --- phụ phẩm nông nghiệp và giấy thải bỏ của công đoạn thêu vật liệu may}. GVHD: TS. Phạm Thị Hồng Phượng.
  \item Môn học nền: Hóa hữu cơ, Hóa lý, Kỹ thuật phản ứng, Quá trình \& thiết bị, Phân tích công cụ, Kỹ thuật an toàn.
\end{cvitems}

\section{Kỹ năng chuyên môn}
\begin{cvfields}[.25\linewidth]
  \cvfield{Đã vận hành}{Máy đo độ nhớt Stormer, Pfund Cryptometer, thước độ mịn Hegman, cốc đo tỷ trọng, máy đo độ bóng, thiết bị đo độ bám dính}
  \cvfield{Đã tiếp cận}{Buồng phun sương muối (ISO 7253), lão hóa nhanh QUV, thử uốn Elcometer, thử va đập; đánh giá gỉ/phồng rộp ISO 4628-2, ASTM D1654}
  \cvfield{Thiết bị phòng TN}{Sử dụng được các thiết bị đo lường kỹ thuật và phân tích cơ bản tại phòng thí nghiệm nhà trường}
  \cvfield{Vật liệu}{Nhựa alkyd, acrylic, epoxy, polyurethane, polyester; hệ dung môi; bột màu --- bột độn; composite từ sợi tự nhiên (lục bình, phụ phẩm nông nghiệp)}
  \cvfield{Hệ thống chất lượng}{ISO 9001:2015, 5S, quy trình IQC/IPQC/FQC, SOP, tra cứu TCVN --- ASTM --- ISO}
  \cvfield{An toàn --- môi trường}{MSDS/SDS, PPE, PCCC, phân loại và xử lý chất thải nguy hại}
  \cvfield{Phần mềm}{ChemDraw, Excel (xử lý số liệu kiểm nghiệm), Word --- soạn báo cáo kỹ thuật}
\end{cvfields}

\section{Nghiên cứu \& Dự án}
\cventry{Khóa luận: Vật liệu composite giả da từ lục bình và phụ phẩm nông nghiệp}{Khoa Công nghệ Hóa học, IUH --- GVHD: TS. Phạm Thị Hồng Phượng}{2026}{TP. Hồ Chí Minh}
\begin{cvitems}
  \item Khảo sát các yếu tố ảnh hưởng đến quy trình chế tạo vật liệu giả da từ \textbf{lục bình, phụ phẩm nông nghiệp và giấy thải} của công đoạn thêu ngành may.
  \item Hướng tới tận dụng phế phụ phẩm, giảm chất thải --- vật liệu thân thiện môi trường thay thế da tổng hợp.
\end{cvitems}

\cventry{Báo cáo thực tập tốt nghiệp (nhóm 8 sinh viên)}{Khoa Công nghệ Hóa học, IUH --- GVHD: ThS. Lê Nhất Thống}{2026}{TP. Hồ Chí Minh}
\begin{cvitems}
  \item Mô tả trọn quy trình sản xuất sơn: nguyên liệu, sơ đồ khối, thiết bị, sự cố thường gặp và biện pháp khắc phục.
  \item Hệ thống hóa \textbf{12 phép thử chất lượng} kèm nguyên lý, cấu tạo thiết bị, tiêu chuẩn áp dụng và quy trình thao tác từng bước.
\end{cvitems}

\section{Giải thưởng}
\begin{cvitems}
  \item \textbf{Giải Nhì \& Giải Ba --- Khoa học Kỹ thuật cấp tỉnh} (lĩnh vực mạng xã hội), bậc THPT.
  \item \textbf{Giải Nhì Cờ vua cấp tỉnh}, bậc THPT.
  \item \textbf{Giải Khuyến khích quốc gia} --- cuộc thi \textit{An toàn giao thông cho nụ cười ngày mai}.
  \item \textbf{Giải Khuyến khích cấp tỉnh} --- cuộc thi Văn hóa đọc.
\end{cvitems}

\section{Chứng chỉ}
\begin{cvfields}[.25\linewidth]
  \cvfield{Tin học}{Chứng chỉ Ứng dụng công nghệ thông tin cơ bản --- Trường Đại học Công nghiệp TP.HCM}
\end{cvfields}

\section{Hoạt động khác}
\cventry{Cộng tác viên}{UBND Phường 5, [Thành phố/Tỉnh]}{[thời gian]}{[địa điểm]}
\begin{cvitems}
  \item Hỗ trợ các hoạt động phong trào và công tác cộng đồng tại địa phương; phối hợp làm việc nhóm, tiếp xúc và trao đổi với người dân.
\end{cvitems}

\end{document}
